\section{Background and Motivation}
\label{sec:background}

\subsection{Slow Faults in Production Systems}

Gray failures~\cite{huang2017gray} encompass a spectrum of partial
failures where components are neither fully operational nor completely
down. Slow faults---a subclass where service time degrades by a factor
$s > 1$---arise from diverse root causes: memory pressure causing
excessive garbage collection, disk degradation increasing I/O latency,
thermal throttling reducing CPU frequency, or network congestion on a
specific path~\cite{gunawi2018fail}.

Unlike crash faults, slow faults are difficult to detect because the
affected node continues to respond, albeit slowly. Standard health
checks (TCP probes, heartbeats) cannot distinguish a 5$\times$ slower
node from a healthy one under transient load. This ambiguity is
especially pronounced at high system utilization, where \emph{all} nodes
exhibit elevated latency due to queueing effects.

\subsection{Power-of-Two-Choices Load Balancing}

Power-of-Two-Choices (P2C)~\cite{mitzenmacher2001power} is the dominant
load balancing algorithm in modern distributed systems, used in
Envoy~\cite{envoy}, nginx, and numerous cloud-native
infrastructure. For each incoming request, P2C samples two candidate
workers at random and routes the request to the one with the shorter
queue.

P2C provides an important but under-appreciated implicit mitigation
mechanism for slow faults. A node with slowdown factor $s$ processes
requests $s$ times slower, causing its queue to grow. P2C's queue-depth
comparison then naturally avoids this node, routing traffic to healthier
alternatives. We show in Section~\ref{sec:nonmono} that this implicit
mitigation is highly effective at extreme severities ($s > \sstar$) but
insufficient at moderate severities ($s < \sstar$), precisely where
explicit mitigation strategies add value.

\subsection{Existing Mitigation Approaches}

\paragraph{Hedged requests.}
Dean and Barroso~\cite{dean2013tail} proposed sending a speculative
copy of a request to a second server after a delay (typically the
expected service time). The client uses whichever response arrives
first. This is effective for random stragglers but, as we demonstrate,
ineffective for persistent slow faults because new requests continue to
be routed to the faulty node.

\paragraph{Outlier blacklisting.}
Systems like Cassandra's dynamic snitch~\cite{lakshman2010cassandra}
track per-node latency and temporarily blacklist nodes whose latency
exceeds $k$ standard deviations above the cluster median. This is a
proactive mechanism that redirects future traffic, but risks oscillation
at high load when blacklisting reduces capacity and increases latency on
remaining nodes.

\paragraph{Circuit breaking.}
Istio and Envoy implement circuit breakers that open when error rates or
latency exceed thresholds, redirecting traffic away from unhealthy
endpoints. These are effective but require manual threshold tuning and
do not adapt to varying load conditions.

\subsection{The Mitigation Trap}
\label{sec:trap}

A key motivation for this work is the \emph{mitigation trap}: at high
system load, naive mitigation strategies can be \emph{worse than doing
nothing}. Consider hedged requests at 85\% load: each hedge adds a
duplicate request, pushing healthy workers from $\rho = 0.85$ toward
saturation. The M/D/1 queueing delay grows as $\rho/(2(1-\rho))$,
creating a positive feedback loop: hedges increase load, which
increases latency, which triggers more hedges.

We discovered this phenomenon empirically (Section~\ref{sec:eval-lit})
and it motivates the load-aware design of our adaptive framework: any
mitigation strategy must account for its own capacity cost.
